\section{Dodatek --- podział pracy}
\label{sec:dodatek}

\noindent
Poniższa tabela prezentuje szczegółowy podział zadań w~projekcie wraz
z~procentowym udziałem członków zespołu. Ze względu na charakter
przedmiotu (Badania Operacyjne) podział został tak skonstruowany, aby
każdy z~członków zespołu uczestniczył w~realnej pracy algorytmicznej
w~porównywalnym wymiarze, przy czym \textbf{rdzeń metaheurystyki ABC}
pozostaje głównym obszarem odpowiedzialności jednego autora.

\begin{longtable}{p{4.4cm} p{5.6cm} p{4.4cm}}
\toprule
\textbf{Etap realizacji} & \textbf{Zakres pracy} & \textbf{Osoba / \% udział} \\
\midrule
\endfirsthead

\toprule
\textbf{Etap realizacji} & \textbf{Zakres pracy} & \textbf{Osoba / \% udział} \\
\midrule
\endhead

Model matematyczny problemu
& Formalizacja jako ILP: zmienne $x_{ij}, y_i$, funkcja celu,
  ograniczenia~\eqref{eq:start}--\eqref{eq:bin}, ograniczenia
  typowe~\eqref{eq:types}
& Adrian 50\%, Mikołaj 30\%, \\[-2pt]
& & Bartosz 20\% \\[4pt]

Algorytm ABC --- rdzeń metaheurystyki
& \texttt{abc\_algorithm.py}: faza employed / onlooker / scout, ruletka
  selekcji proporcjonalnej, licznik prób, reset rozwiązań,
  pętla główna i~kryterium stopu
& \textbf{Bartosz 70\%}, \\[-2pt]
& & Adrian 15\%, Mikołaj 15\% \\[4pt]

Operatory zaburzeń
& Operatory mutacji rozwiązania (add / remove / swap / 2-opt / shift)
  wraz z~rozkładem prawdopodobieństw oraz mechanizm utrzymywania
  dopuszczalności po~mutacji
& \textbf{Bartosz 65\%}, \\[-2pt]
& & Mikołaj 20\%, Adrian 15\% \\[4pt]

A\textsuperscript{*} i konstrukcja ścieżek
& Implementacja A\textsuperscript{*} na siatce 8-kierunkowej, twardy
  zakaz powtarzania krawędzi, dynamiczne blokowanie atrakcji
  spoza listy waypointów, sklejanie segmentów trasy
& Adrian 55\%, Bartosz 30\%, \\[-2pt]
& & Mikołaj 15\% \\[4pt]

Funkcja ewaluacji
& \texttt{evaluate\_path}: wartość trasy, koszt ruchu, koszt atrakcji,
  czas; leksykograficzne kryterium porównawcze rozwiązań
& Adrian 60\%, Bartosz 25\%, \\[-2pt]
& & Mikołaj 15\% \\[4pt]

Generator mapy
& \texttt{map\_generator.py}, \texttt{src/map\_generator.py}:
  procedualne generowanie wag terenu, rozmieszczanie atrakcji,
  przypisanie typów, parametry rozkładu
& \textbf{Mikołaj 65\%}, \\[-2pt]
& & Adrian 20\%, Bartosz 15\% \\[4pt]

Populacja początkowa
& \texttt{initial\_population.py}: heurystyczna konstrukcja rozwiązań
  dopuszczalnych, dobór punktów pośrednich z~uwzględnieniem
  ograniczeń typowych, kontrola czasu/budżetu
& Mikołaj 45\%, Bartosz 40\%, \\[-2pt]
& & Adrian 15\% \\[4pt]

Walidator ograniczeń
& \texttt{examples/validate\_population.py}: sprawdzanie
  spełnienia~\eqref{eq:start}--\eqref{eq:bin}, weryfikacja braku
  powtórzeń krawędzi i~ograniczeń typowych
& Adrian 65\%, Mikołaj 25\%, \\[-2pt]
& & Bartosz 10\% \\[4pt]

Parser i I/O
& \texttt{parser.py}, \texttt{load\_json.py}: konwersja JSON $\leftrightarrow$
  struktury wewnętrzne, walidacja schematu wejścia
& Mikołaj 55\%, Adrian 35\%, \\[-2pt]
& & Bartosz 10\% \\[4pt]

Wizualizacja tras
& \texttt{visualize\_paths.py}: rysowanie tras, atrakcji, punktów
  start/end, kolorowanie typów, wykresy diagnostyczne
& Mikołaj 70\%, Adrian 20\%, \\[-2pt]
& & Bartosz 10\% \\[4pt]

Aplikacja web (Streamlit)
& \texttt{app.py}, \texttt{src/path\_solver.py}, \texttt{src/ui\_io.py},
  \texttt{src/ui\_viz.py}: edytor mapy, generator, walidacja,
  tryb wyników
& Mikołaj 50\%, Adrian 35\%, \\[-2pt]
& & Bartosz 15\% \\[4pt]

Scenariusze testowe
& \texttt{scenarios/small.json}, \texttt{scenarios/medium.json}:
  parametry, dobór trudności, dane wejściowe do eksperymentów
& Adrian 45\%, Mikołaj 35\%, \\[-2pt]
& & Bartosz 20\% \\[4pt]

Dokumentacja LaTeX
& Niniejszy dokument: opis problemu, model matematyczny, opis
  algorytmu, dokumentacja użytkownika, literatura
& Adrian 40\%, Mikołaj 35\%, \\[-2pt]
& & Bartosz 25\% \\

\bottomrule
\end{longtable}

\subsection*{Sumaryczny rozkład wkładu}

Po~zsumowaniu wkładów z~poszczególnych etapów (z~wagami proporcjonalnymi
do nakładu pracy w~danym module) orientacyjny rozkład wkładu autorów
w~całość projektu przedstawia się następująco:

\begin{center}
\begin{tabular}{l c}
\toprule
\textbf{Autor} & \textbf{Łączny wkład} \\
\midrule
Bartosz Tokarz   & $\approx$ 34\% \\
Adrian Michalski & $\approx$ 34\% \\
Mikołaj Kaleta   & $\approx$ 32\% \\
\bottomrule
\end{tabular}
\end{center}

\noindent
\textbf{Komentarz do podziału.}
Bartosz Tokarz odpowiada za~rdzeń metaheurystyki ABC --- najbardziej
wymagający koncepcyjnie i~obliczeniowo komponent projektu (główna
pętla, operatory zaburzeń, mechanizm scout). Adrian Michalski skupia
się na~komponentach algorytmicznych powiązanych ze~ścieżkami
i~poprawnością rozwiązań: A\textsuperscript{*} z~ograniczeniami,
funkcja ewaluacji, walidator i~formalizacja modelu matematycznego.
Mikołaj Kaleta odpowiada za~algorytmikę po~stronie środowiska:
procedualny generator mapy, konstrukcja populacji początkowej oraz
warstwa wizualizacji i~I/O. Taki podział utrzymuje porównywalny wkład
algorytmiczny wszystkich członków zespołu, przy jednoznacznym
przypisaniu odpowiedzialności za~główny algorytm projektu.
